\documentclass[11pt,a4paper]{article}
\usepackage[margin=1in]{geometry}
\usepackage{amsmath,amssymb}
\usepackage{parskip}
\usepackage{hyperref}

\title{\textbf{Pre-Experiment Overview:\\ Elementary Charge and Millikan Experiment}}
\author{IBRAHIM H.I. ABUSHAWISH \\ Istanbul University, Department of Physics}
\date{}

\begin{document}
\maketitle

\section*{Purpose of the Experiment}
The main goal of this lab is to observe tiny, charged oil droplets as they are accelerated between two capacitor plates. You will be timing exactly how fast these droplets rise and fall when subjected to different voltages. By measuring these velocities---both moving with gravity and fighting against it---you can calculate the specific charge and radius of each individual droplet. Ultimately, you use these measurements to determine the fundamental elementary charge and the $e/m$ ratio of the electron.

\section*{Importance and Significance}
The real beauty of this classic Millikan setup is that it physically demonstrates that electric charge is quantized. The entire experiment relies on balancing several competing forces acting on a single droplet: gravitational pull, air buoyancy, the applied electric field, and viscous drag, which is dictated by Stokes' Law.

By adding and subtracting the equations for a droplet's rising and falling velocities, you can isolate its total charge. The magic happens when you plot the calculated charges of multiple droplets against their radii: the data points don't scatter randomly, but instead align into distinct horizontal levels. The gap between these levels proves that a droplet's total charge is always a strict integer multiple of one base value, expressed as $Q = n \cdot e$. By measuring the distance between these layered lines on your graph, you are directly determining the elementary charge itself, which the literature sets at $e = 1.602 \times 10^{-19}\text{ C}$. You are essentially catching the fundamental building block of electricity in plain sight!

\end{document}
